\chapter{Versuchsmethodik}
\label{ch:strategy}

Dieses Kapitel beschreibt den geplanten Versuch. Es legt fest, wie die
Balken beschafft und ausgewaehlt werden sollen, wie die Proben aus dem
Querschnitt entnommen werden, mit welchen Pruefverfahren die
\gls{restfestigkeit} bestimmt wird und wie die Messwerte ausgewertet werden
sollen. Der Versuch ist zum Zeitpunkt der Abfassung dieses Kapitels noch
nicht durchgefuehrt. Alle Angaben sind daher als Planung zu lesen, die vor
Versuchsbeginn mit der Betreuung abzustimmen und nach einem Vorversuch
nachzuschaerfen ist.

\section{Herkunft und Auswahl der Balken}
\label{sec:material-herkunft}

Vorgesehen ist die Verwendung von \gls{altholz} aus einem Rueckbau, dessen
Einbausituation noch zugaenglich ist. Nur so laesst sich die \gls{lastgeschichte}
des Balkens ueberhaupt rekonstruieren, und nur dann ist der Vergleich zwischen
frueherer \gls{druckzone} und frueherer \gls{zugzone} inhaltlich tragfaehig.
Fuer jeden Balken soll vor oder waehrend des Ausbaus dokumentiert werden,
welche Seite im Tragwerk oben lag, an welcher Position im Tragwerk der Balken
lag, welche Spannweite und welche Auflagerausbildung vorlagen und welche
Nutzung ueber dem Balken stattfand. Zusaetzlich ist eine fotografische
Aufnahme der Oberseite sowie eine dauerhafte Markierung dieser Seite
vorzunehmen, die ein spaeteres Zersaegen ueberlebt, etwa eine Punzung oder
eine eingebrachte Gravur mit Datum und Herkunft. Wo es die Zugaenglichkeit
erlaubt, soll die Durchbiegung des Balkens vor dem Rueckbau gemessen werden,
um das ueber Jahrzehnte akkumulierte \gls{kriechen} als beobachtete Groesse
festzuhalten und nicht lediglich zu unterstellen.

Diese Dokumentationsschritte sind in keiner der einschlaegigen Normen
geregelt. Sie sind als Best Practice zu verstehen, nicht als normative
Vorgabe, und genau darin liegt ein methodisches Risiko: die Guete der
Zuordnung haengt am Sorgfaltsgrad der Erhebung und laesst sich nachtraeglich
nicht mehr herstellen.

\section{Das Identifikationsproblem der Einbaulage}
\label{sec:identifikation}

Der kritische Punkt der gesamten Untersuchung liegt vor der ersten Pruefung.
Nach dem Ausbau aus dem Tragwerk ist nicht mehr unmittelbar sichtbar, welche
Seite des Querschnitts unter Druck und welche unter Zug stand. Ein Balken,
der ohne Zuordnung auf dem Lagerplatz liegt, traegt die fuer diese Arbeit
entscheidende Information nicht mehr an sich. Wird die Orientierung
vertauscht, so kehrt sich das Vorzeichen der gemessenen Differenz um, und
mehrere vertauschte Balken loeschen einen real vorhandenen Effekt im Mittel
aus. Das Problem ist damit nicht nur eine Unschaerfe, sondern ein Fehler, der
systematisch gegen die Hypothese arbeitet.

Liegt keine Dokumentation vor, so kommen nur Sekundaerindikatoren in Betracht.
Denkbar sind Nagelspuren oder Auflagerabdruecke von Dielen auf der ehemaligen
Oberseite, eine seitenweise unterschiedliche Ausbildung von Oberflaechenrissen
sowie Feuchteverfaerbungen oder Pilzbefall, die bei langfristigem Feuchtestau
unterseitig ausgepraegter auftreten koennen. Auch der radiale Abstand zum Kern
kann einen Hinweis geben. Alle diese Indikatoren sind unsicher und in der
Literatur nicht als Verfahren abgesichert; sie werden hier ausdruecklich als
Plausibilitaetshinweise und nicht als Nachweis behandelt.

Daraus folgt ein hartes Ausschlusskriterium: Balken, deren Herkunft unbekannt
oder widerspruechlich dokumentiert ist und bei denen die Sekundaerindikatoren
keine eindeutige Aussage liefern, werden nicht geprueft. Im Zweifelsfall wird
die Probe ausgeschlossen. Eine geringere Zahl sicher zugeordneter Balken ist
einer groesseren Zahl zweifelhaft zugeordneter vorzuziehen, weil letztere den
Effekt verwaessern, ohne die Aussagekraft zu erhoehen. Jeder Ausschluss wird
mit Grund und Datum im Pruefprotokoll vermerkt, damit die Ausschlussquote
selbst ein berichtetes Ergebnis der Arbeit bleibt.

\section{Probenentnahme aus Druck- und Zugzone}
\label{sec:probenentnahme}

Aus jedem zugelassenen Balken soll ein Segment von etwa \SI{300}{\milli\metre}
bis \SI{500}{\milli\metre} Laenge herausgesaegt werden, und zwar ausserhalb der
Bereiche mit konstruktiv bedingter Fehlerkonzentration, also nicht an
Auflagern und nicht an Anschluessen. Dieses Segment wird entlang der
horizontalen Mittellinie geteilt, die als Lage der \gls{neutralefaser}
angenommen wird. Die Oberhaelfte gilt als vermutete \gls{druckzone}, die
Unterhaelfte als vermutete \gls{zugzone}. Die Proben werden in beiden Haelften
mit gleichem Abstand zur Trennebene entnommen, damit der Abstand zur neutralen
Faser fuer die beiden verglichenen Zonen derselbe ist und die randnahen,
ehemals am hoechsten beanspruchten Fasern in beiden Zonen gleichermassen
erfasst werden. Randnahe und kernnahe Lage werden nicht gemischt, weil sonst
die Position im Querschnitt mit der Zonenzugehoerigkeit verwechselt wuerde.

Jede Probe erhaelt eine Kennung aus Balken-Identifikation, Zone und Laufnummer,
die durch das Saegen nicht verloren gehen kann, etwa durch Punzung. Je Halbsegment
sind mindestens drei Proben pro Pruefverfahren vorgesehen, angestrebt werden
vier bis fuenf, sodass pro Balken mit einer Groessenordnung von 24 bis 40
Kleinproben zu rechnen ist.

\section{Konditionierung, Holzfeuchte und Rohdichte}
\label{sec:konditionierung}

Nach Ausbau und Transport sollen die Balkensegmente zunaechst trocken,
beschattet und bei moderatem Klima gelagert werden, bevor sie zerteilt werden;
vorgesehen ist ein Zeitraum von mehreren Wochen, um oberflaechennahe
Trocknungsspannungen abzubauen. Diese Wartezeit ist eine Vorsichtsmassnahme
ohne normative Grundlage. Die fertig gesaegten Kleinproben werden anschliessend
im Normalklima von \SI{20}{\celsius} und \SI{65}{\percent} relativer
Luftfeuchte konditioniert, bis sich die Holzfeuchte auf einem stabilen Niveau
eingestellt hat.

Die Holzfeuchte wird nach der Darrmethode bestimmt
\parencite{din13183}: Waegung der feuchten Probe, Trocknung bei
\SI{103}{\celsius} bis zur Massenkonstanz, Waegung der Trockenmasse und
Berechnung der Feuchte als Massenverlust bezogen auf die Trockenmasse. Die
\gls{rohdichte} wird an Staebchen aus dem Bruchstueck geometrisch bestimmt,
also aus Masse und aus den auf \SI{0.1}{\milli\metre} gemessenen Kantenlaengen
\parencite{iso13061}, und nach der in \textcite{en384} vorgesehenen Systematik
auf die Bezugsfeuchte von \SI{12}{\percent} umgerechnet. Holzfeuchte und
Rohdichte sind fuer diese Arbeit keine Nebengroessen: beide beeinflussen die
Festigkeit unmittelbar und muessen daher je Zone einzeln vorliegen, damit ein
gemessener Unterschied nicht faelschlich der Lastgeschichte zugeschrieben wird.

\section{Pruefverfahren: Kleinproben statt Bauteilpruefung}
\label{sec:pruefverfahren}

Zwei fachliche Klarstellungen bestimmen die Wahl der Verfahren.

Erstens liefert der Vierpunkt-Biegeversuch die Biegefestigkeit
(\acrshort{mor}) und den Elastizitaetsmodul (\acrshort{moe}), nicht die
Querzugfestigkeit \parencite{en408,fourpoint}. Die Festigkeit quer zur Faser
waere in einem eigenen Zugversuch quer zur Faser an eigener Probengeometrie zu
bestimmen. Die Biegefestigkeit ist hier die Hauptzielgroesse; eine Bestimmung
der Querzugfestigkeit bleibt optional und ist nicht Teil des geplanten
Versuchsumfangs.

Zweitens ist \textcite{en408} fuer Proben in Bauteilabmessungen ausgelegt. Da
Druck- und Zugzone desselben Balkens verglichen werden sollen, muss der
Querschnitt geteilt werden, und die entstehenden Proben sind Kleinproben ohne
sichtbare Fehler. Die normative Grundlage wechselt daher auf die
Kleinprobennormen: DIN 52186 beziehungsweise ISO 13061-3 fuer die
Biegefestigkeit parallel zur Faser, \textcite{din52185} beziehungsweise
ISO 13061-17 fuer die Druckfestigkeit parallel zur Faser
\parencite{iso13061} sowie die oben genannten Normen fuer Dichte und
Holzfeuchte. Die Geometrieprinzipien aus \textcite{en408} werden sinngemaess
uebernommen, namentlich eine Stuetzweite von 18 mal der Querschnittshoehe und
ein Abstand der inneren Lasteintragungspunkte von 6 mal der Hoehe. Dies ist
der methodische Kern des Kapitels: nicht EN 408 unmittelbar, sondern
Kleinprobennormen mit den Geometrieprinzipien von EN 408.

Zu den Querschnittsabmessungen der Kleinproben, den Belastungsgeschwindigkeiten
und der Mindestlaenge der Druckproben liegen in der ausgewerteten Quellenlage
keine gesicherten Festlegungen vor; sie sind vor Versuchsbeginn anhand der
Normtexte selbst und der verfuegbaren Balkenstaerke festzulegen. Waehrend der
Pruefung werden Bruchlast, Bruchbild und Bruchlage protokolliert. Ein Bruch
ausserhalb des Bereichs zwischen den inneren Lastpunkten, ein Schubbruch
anstelle eines Faserbruchs sowie ein Stabilitaetsversagen durch Knicken im
Druckversuch fuehren zum Ausschluss des jeweiligen Messwerts, weil in diesen
Faellen nicht die gesuchte Materialeigenschaft gemessen wurde. Aeste, Risse,
Faeule und starke Faserabweichung im Prueflaengenbereich fuehren bereits vor
der Pruefung zum Ausschluss, entsprechend der Logik der visuellen
\gls{sortierung} nach \textcite{din4074}.

\section{Statistische Auswertung}
\label{sec:auswertung}

Druck- und Zugprobe stammen aus demselben Balkensegment an derselben
Laengsposition. Sie sind daher nicht unabhaengig, sondern gepaart, und der
Vergleich erfolgt als gepaarter Test und nicht als Vergleich zweier
unabhaengiger Gruppen. Fuer jedes Segment wird die Differenz der Zielgroesse
zwischen Zugzone und Druckzone gebildet, und die Auswertung findet auf diesen
Differenzen statt. Dieser Aufbau ist der eigentliche Vorzug des Versuchs: die
grosse Streuung zwischen verschiedenen Balken, die Holzart, Wuchsort und Alter
mitbringen, faellt bei der Differenzbildung heraus.

Die Verteilungsannahme wird vor dem Haupttest geprueft, mit einem
Shapiro-Wilk-Test auf die Differenzen und mit einem Q-Q-Diagramm zur
visuellen Kontrolle; die Varianzhomogenitaet zwischen den Zonen wird mit einem
Levene-Test betrachtet. Sind die Differenzen annaehernd normalverteilt, wird
zweiseitig ein gepaarter t-Test auf dem Signifikanzniveau von
\SI{5}{\percent} gerechnet, andernfalls der Wilcoxon-Vorzeichen-Rang-Test als
verteilungsfreie Alternative. Berichtet werden neben dem p-Wert die
Effektgroesse und das 95-Prozent-Konfidenzintervall der mittleren Differenz,
da eine reine Signifikanzaussage bei kleinen Stichproben wenig aussagt.

Die Rohdichte wird als Kovariate gefuehrt. Zunaechst wird geprueft, ob sich die
Rohdichte zwischen den beiden Zonen systematisch unterscheidet. Ist das der
Fall, wird der Festigkeitsvergleich um die Dichte bereinigt, im Rahmen einer
Kovarianzanalyse als Erweiterung der \acrshort{anova}. Andernfalls koennte ein
Dichteunterschied innerhalb des Querschnitts als Wirkung der Lastgeschichte
missdeutet werden. Werden mehrere Segmente je Balken entnommen, so liegt eine
genestete Struktur vor, und die Auswertung erfolgt in einem gemischten Modell
mit der Zone als festem Effekt und dem Balken als Zufallseffekt, um
Pseudoreplikation zu vermeiden.

\section{Stichprobenumfang}
\label{sec:stichprobenumfang}

Als Planungsgroesse werden zwoelf bis fuenfzehn Balkensegmente angesetzt, wie
sie in vergleichbaren Altholzuntersuchungen ueblich sind, mit vier Proben je
Zone fuer Biegung und Druck
sowie zwei Proben je Zone fuer Holzfeuchte und Rohdichte
(\Cref{tab:versuchsplan}). Diese Zahl ist ausdruecklich als konservative
Untergrenze zu lesen und nicht als normative Vorgabe: weder
\textcite{en384} noch \textcite{din52185} legen fuer den Zonenvergleich an
Altholz eine Mindestzahl fest, und die genannte Groessenordnung ist
erfahrungsgestuetzt, nicht normativ abgesichert. Vorgesehen ist daher ein Vorversuch an wenigen Segmenten, aus dessen
Effektgroesse eine Poweranalyse den erforderlichen Umfang fuer eine
angestrebte Teststaerke von \SI{80}{\percent} bei einem Signifikanzniveau von
\SI{5}{\percent} neu berechnet. Faellt der beobachtete Effekt klein aus, so
kann der erforderliche Umfang die verfuegbare Menge sicher zugeordneter Balken
deutlich uebersteigen. In diesem Fall ist dies als Ergebnis zu berichten und
nicht durch Aufnahme unsicher zugeordneter Balken zu umgehen. Einzuplanen ist
ferner eine Ausfallquote durch Ausschlusskriterien, Messfehler und ungueltige
Bruchbilder.

\begin{table}[H]
\centering
\caption{Geplanter Versuchsumfang als konservative Untergrenze, vor dem
Vorversuch und der darauf gestuetzten Poweranalyse.}
\label{tab:versuchsplan}
\begin{tabularx}{\textwidth}{@{}Xccc@{}}
\toprule
\textbf{Pruefverfahren} & \textbf{Proben je Zone} & \textbf{Segmente} &
\textbf{Proben gesamt} \\
\midrule
Vierpunkt-Biegung (\acrshort{mor}, \acrshort{moe}) & 4 & 12 bis 15 & 96 bis 120 \\
Druckversuch parallel zur Faser & 4 & 12 bis 15 & 96 bis 120 \\
Holzfeuchte (Darrmethode) & 2 & 12 bis 15 & 48 bis 60 \\
Rohdichtebestimmung & 2 & 12 bis 15 & 48 bis 60 \\
\midrule
\textbf{Summe} & & \textbf{12 bis 15} & \textbf{288 bis 360} \\
\bottomrule
\end{tabularx}
\end{table}

\section{Grenzen der Messung}
\label{sec:grenzen-messung}

Abschliessend ist der Anspruch der Messung zu begrenzen. Der Versuch bestimmt
die Festigkeit der beiden Zonen im Verhaeltnis zueinander, nicht einen
absoluten Verlust gegenueber Neuholz. Ein Bezug auf Festigkeitsklassen nach
\textcite{en338} oder auf die aelteren Untersuchungen an altem
Konstruktionsholz \parencite{rug1989,ehlbeck1990} bleibt daher ein
Literaturvergleich und kein eigener Nachweis.
